%!TEX root = ../../../main.tex

\chapter{Problem, Idee und Daten}
\label{chap:problem}

\paragraph{Problem.} Zu Hause sammeln sich Rechnungen und Verträge. Will ich etwas wissen, muss ich die richtige Stelle in der richtigen Datei finden, oft in zwei Dateien (\enquote{Stimmt der Abschlag der Rechnung mit dem Vertrag überein?}). Mein Agent beantwortet solche Fragen über alle Dokumente, belegt jede Aussage mit einer nummerierten Quelle und sagt, wenn nichts belegt ist. Weil die Dokumente IBAN, Adressen und Gesundheitsdaten enthalten, läuft alles lokal auf meinem Laptop (16\,GB RAM, keine Grafikkarte) mit Ollama und ChromaDB.

\paragraph{Fragetypen.} Meine eigenen Fragen fallen in fünf Typen (Tabelle~\ref{tab:fragetypen}). Sie bestimmen die Architektur: Für a und b reicht eine gute Suche, c braucht zwei Dokumente zugleich, d alle Verträge, und e muss der Agent als unbeantwortbar erkennen.

\begin{table}[H]
    \centering\footnotesize\setstretch{1.05}
    \begin{tabularx}{\textwidth}{@{}l Y Y@{}}
        \toprule
        \textbf{Typ} & \textbf{Beispiel} & \textbf{Was der Agent dafür braucht} \\
        \midrule
        a) Wert aus Tabelle & Wie viel Umsatzsteuer weist die LuminaStyle-Rechnung aus? & die richtige Tabellenzeile mit Spaltenkopf \\
        b) Textstelle & Wie lange läuft der Handyvertrag mindestens? & den richtigen Abschnitt im Vertrag \\
        c) Zwei Dokumente & Passt die Vorauszahlung der Wasserabrechnung zum Vertrag? & Rechnung \emph{und} Vertrag, die zusammengehören \\
        d) Liste & Welche Verträge verlängern sich automatisch um ein Jahr? & alle Verträge, Fristen richtig berechnet \\
        e) ohne Beleg & Welche Kündigungsfrist hat mein Handyvertrag bei Vodafone? & erkennen, dass es kein Dokument gibt \\
        \bottomrule
    \end{tabularx}
    \caption{Fragetypen und was sie vom Agenten verlangen}
    \label{tab:fragetypen}
\end{table}

\paragraph{Eigener Datensatz.} Die Beispielfragen der Vorlesung (Prüfungsordnung, Modulhandbuch) passen nicht zu meinem Use Case. Ich habe daher einen eigenen Datensatz gebaut: zehn Rechnungen als PDF mit echten Layouts, darunter Tabellen mit Hoch- und Niedertarif und ein Diagramm, dessen Werte nur im Bild stehen, und sechs Verträge (drei DOCX, drei PDF). Die Verträge erzeugt ein Skript aus Sollwerten, so kenne ich zu jeder Frage Antwort und Dokument. Zwei Fehler sind Absicht: Der Heizstromvertrag nennt 225\,€ Abschlag, die Rechnung 220\,€, und der Wasservertrag ist nicht kündbar. Daran sehe ich, ob der Agent Unterschiede bemerkt.
